\section{Conclusion}

The use of gadgets as a learning medium is a topic that can be analyzed
using various statistical concepts. Based on data from 69 respondents,
the analysis produced a Pearson correlation coefficient of 0.733,
indicating a strong positive relationship between variables X and Y.

The analysis also produced an $R^2$ value of 53.66\%. This means that,
within the research model, approximately 53.66\% of the variation in
variable Y can be explained by variable X, while the remaining 46.34\% is
not explained by X within the model and may be associated with other
factors.

The t-test produced $t = 8.808$ with a significance value of 0.000,
indicating that variable X has a statistically significant relationship
with variable Y partially within the tested model. Meanwhile, the F-test
produced $F = 77.576$ with a significance value of 0.000, indicating that
the regression model as a whole is statistically significant.

In addition, the concept of probability can be directly applied to the 69
respondents. Using $\ge 60$ as the threshold for the high category, we
obtained a probability of 43.48\% for high X, 13.04\% for high Y, and
10.14\% for having both high X and high Y. When it is known that a
respondent has a high X score, the probability of having a high Y score
becomes 23.33\% through conditional probability.

This learning journey shows that different statistical concepts are
connected. We started by understanding the data, summarizing it through
descriptive statistics, identifying patterns through visualization,
analyzing relationships using correlation and regression, testing
statistical significance through the t-test and F-test, and finally
applying probability to understand the likelihood of certain conditions
occurring.

Ultimately, we learned that the purpose of statistics is not simply to
produce tables, graphs, and numbers, but to transform data into
meaningful insights that can be understood, explained logically, and
interpreted carefully within the context of the research.
